\documentclass[../../../include/open-logic-section]{subfiles}

\begin{document}

\olfileid{sth}{replacement}{absinf}
\olsection[प्रतिस्थापन आणि निरपेक्ष अनंतता]{प्रतिस्थापन आणि ``निरपेक्ष अनंतता''}

आता \limofsize{} बाजूला ठेवून प्रतिस्थापनाचे समर्थन करण्याचे वेगळ्या प्रकारचे (आंतरिक) प्रयत्न पाहू. हे प्रयत्न संच टप्प्याटप्प्याने तयार होतात या कल्पनेला गांभीर्याने घेतात.

क्रमिक प्रक्रिया प्रथम सांगताना आपण प्रत्येक टप्प्यावर काय घडते हे स्पष्ट करणारी काही तत्त्वे मांडली होती: \stageshier, \stagesord आणि \stagesacc. नंतर टप्प्यांच्या संख्येविषयी काही सांगणारी तत्त्वे जोडली: \stagessucc{} नुसार संच तयार होण्याची प्रक्रिया कधीही संपत नाही, आणि \stagesinf{} नुसार ती अनंताव्या टप्प्यापर्यंत पोहोचते.

ही नंतरची दोन तत्त्वे आणखी व्यापक तत्त्वातून निष्पन्न होतात असे सुचवणे रास्त आहे; उदाहरणार्थ:
\begin{enumerate}
  \item[] \stagesinex. टप्प्यांची संख्या निरपेक्ष अर्थाने अनंत आहे; पदानुक्रम जितका उंच असू शकतो तितका उंच आहे.
\end{enumerate}
अर्थात हे एक अनौपचारिक तत्त्व आहे. पण संचांच्या संचयी-क्रमिक संकल्पनेतून ते लगेच \emph{निष्पन्न} होत नसले, तरी त्या संकल्पनेशी निश्चितच \emph{सुसंगत} वाटते. निदान \limofsize{} च्या विपरीत, संच टप्प्याटप्प्याने तयार होतात ही कल्पना त्यात टिकून राहते.

आता \stagesinex{} वापरून प्रतिस्थापनाचे समर्थन करता येईल अशी आशा आहे. ते कसे शक्य होईल ते पाहू.

\olref[ordinals][idea]{sec} मध्ये आपण पाहिले की (अनौपचारिक अर्थाने) $\omega+\omega$ शी समरूप असायला हवे असे सुक्रमण सहज तयार करता येते. दुसऱ्या शब्दांत, (अनौपचारिक अर्थाने) $\omega+\omega$ हा क्रमप्रकार असलेली टप्प्याटप्प्यांची क्रमिक प्रक्रिया आपण सहज कल्पू शकतो. म्हणून \stagesinex{} स्वीकारले असेल, तर पदानुक्रमात किमान $\omega+\omega$-वा टप्पा, म्हणजे $V_{\omega+\omega}$, आहे हेही स्वीकारले पाहिजे; कारण पदानुक्रम इतक्यापर्यंत \emph{नक्कीच जाऊ शकतो}.

हा विचार सर्वसाधारणपणे असा मांडता येतो: कोणतेही सुक्रमण घ्या; क्रमिक पदानुक्रम बांधण्याची प्रक्रिया किमान त्या सुक्रमणाइतकी पुढे गेली पाहिजे. \olref[ordinals][ordtype]{thmOrdinalRepresentation} याला \emph{स्वयंसिद्धक} मानल्यास याची खात्री मिळेल. मग प्रत्येक सुक्रमण एखाद्या फोन न्यूमन क्रमसूचक संख्येशी समरूप आहे असे ठरेल. प्रत्येक फोन न्यूमन क्रमसूचक संख्येचा प्रतांक ती स्वतःच असल्याने, \olref[ordinals][ordtype]{thmOrdinalRepresentation} यावरून असे ठरेल की संच उपपत्तीत आपण कोणतेही सुक्रमण वर्णन करू शकलो, तर संच बांधण्याची क्रमिक प्रक्रिया त्या सुक्रमणाच्याही पुढे गेली पाहिजे.

हा विचार \stagesinex{} चा परिणाम वाटतो. दुर्दैवाने, त्यावरून प्रतिस्थापन मिळवायचे असेल, तर एक सोपा तांत्रिक अडथळा आहे: प्रतिस्थापन हे \olref[ordinals][ordtype]{thmOrdinalRepresentation} पेक्षा काटेकोरपणे अधिक सामर्थ्यवान आहे. (ही बाब \citet[\S13.2]{Potter2004} यांनी नोंदवली आहे; तिचा पुरावा आपण \olref[sth][replacement][finiteaxiomatizability]{sec} मध्ये देऊ.)

म्हणून \stagesinex{} चा अर्थ प्रतिस्थापन मिळेल अशा रीतीने घ्यायचा असेल, तर पदानुक्रम प्रत्येक सुक्रमणाच्या पुढे जातो \emph{एवढेच} त्याने सांगून चालणार नाही. त्याहून बलवान दावा आवश्यक आहे. यासाठी \cite{Shoenfield:AST} यांनी या कल्पनेचा एक अत्यंत स्वाभाविक विस्तार सुचवला: पदानुक्रम कोणत्याही संचाशी \emph{सहअंतिम} नाही.\footnote{G\"odel यांनीही असाच विचार सुचवला असावा; \citet[p.~223]{Potter2004} पाहा. G\"odel आणि \citeauthor{Shoenfield:AST} यांच्या चर्चेसाठी \citet[90--5]{Incurvati2020} पाहा.} अधिक स्पष्टपणे: $\tau$ हे संचांना पदानुक्रमातील टप्प्यांकडे पाठवणारे प्रतिचित्रण असेल, तर कोणत्याही संच $A$ ची $\tau$ खालील प्रतिमा संपूर्ण पदानुक्रम संपवू शकत नाही. दुसऱ्या शब्दांत (रूपबंधाच्या स्वरूपात):
\begin{enumerate}
  \item[] \stagescofin. $A$ हा संच असेल आणि प्रत्येक $x \in A$ साठी $\tau(x)$ हा टप्पा असेल, तर प्रत्येक $x \in A$ च्या $\tau(x)$ नंतर येणारा एक टप्पा असतो.
\end{enumerate}
$\ZF$ मध्ये \stagescofin{} चे योग्य औपचारिक रूप सिद्ध होते, हे स्पष्ट आहे. उलट, \stagescofin{} प्रतिस्थापनाचे समर्थन करते, असा अनौपचारिक युक्तिवाद करता येतो.\footnote{\stagescofin{} च्या एखाद्या औपचारिक रूपावरून प्रतिस्थापन सिद्ध करणे अधिक कठीण ठरेल; कारण केवळ $\Z$ मध्ये टप्प्यांची व्याख्या करण्याइतके सामर्थ्य नाही. त्यामुळे \stagescofin{} चे औपचारिकीकरण कसे करावे हे स्पष्ट नाही. मात्र \olref[sth][replacement][extrinsic]{sec} मध्ये चर्चिल्याप्रमाणे $\LT$ च्या एखाद्या विस्तारात काम करणे हा एक पर्याय आहे.} समजा $(\forall x \in A)\lexists![y][\phi(x,y)]$. मग प्रत्येक $x \in A$ साठी $\phi(x,y)$ पूर्ण करणारा $y$ म्हणजे $\sigma(x)$ मानू, आणि $\sigma(x)$ पहिल्यांदा ज्या टप्प्यावर तयार होतो तो $\tau(x)$ मानू. \stagescofin{} नुसार असा टप्पा $V$ आहे की $(\forall x \in A)\tau(x)\in V$. आता प्रत्येक $\tau(x) \in V$ आणि $\sigma(x) \subseteq \tau(x)$ असल्याने, विभक्तीकरण वापरून $\Setabs{y
\in V}{(\exists x \in A)\sigma(x) = y} = \Setabs{y}{(\exists x \in
A)\phi(x,y)}$ हा संच मिळवता येतो.

\begin{prob}
  $\ZF$ मध्ये \stagescofin{} चे औपचारिकीकरण करा.
\end{prob}

म्हणून \stagescofin{} प्रतिस्थापनाला समर्थन देते. आणि \stagesinex{} हे \stagescofin{} ला समर्थन देते, हे निदान संभवनीय वाटते. समजा \stagescofin{} खोटे ठरते. म्हणजे पदानुक्रम एखाद्या संच~$A$ शी सहअंतिम आहे: असे $\tau$ प्रतिचित्रण आहे की प्रत्येक टप्पा~$S$ साठी असा काही $x \in A$ आहे की $S \in \tau(x)$. अशा परिस्थितीत पदानुक्रमाची कथित ``निरपेक्ष अनंतता'' पकडण्याचा मार्ग आपल्याकडे आहे: $A$ वर लागू केलेल्या $\tau$ च्या परासाने ती \emph{संपूर्णपणे व्यापली} जाते. त्यामुळे पदानुक्रम ``निरपेक्षपणे अनंत'' आहे या विचाराला बाधा येते. व्यत्यासाने: \stagesinex{} वरून \stagescofin{} मिळते; आणि ते प्रतिस्थापनाचे समर्थन करते.

प्रतिस्थापनाचे आंतरिक समर्थन देण्याचा हा खरोखर आशादायक प्रयत्न आहे. मात्र तो अखेरीस सफल होतो की नाही, हे ठरवण्याचे काम आम्ही तुम्हालाच सोपवतो.

\end{document}
